\section{Methodology}
\label{sec:method}

Our goal is to measure how strongly each pretraining paradigm encodes spurious attributes in frozen representations, holding backbone architecture constant. The methodology is designed around a specific concern: label correlation in supervised training could confound naive comparisons, so every design choice is motivated by the need to isolate the pretraining objective effect.

\subsection{Pretraining Paradigms}

We compare four ResNet-50 backbones pretrained on ImageNet-1k under different objectives:

\begin{itemize}
\item \textbf{Supervised ERM}: Standard cross-entropy classification on ImageNet-1k labels (torchvision \texttt{resnet50(pretrained=True)}).
\item \textbf{MoCo-v3} (contrastive): Momentum-contrast contrastive SSL \citep{He2020MoCo}; official \texttt{r-50-1000ep.pth.tar} checkpoint loaded directly.
\item \textbf{DINO} (self-distillation): Self-distillation with knowledge distillation from a momentum teacher \citep{Caron2021DINO}; \texttt{dino\_resnet50} via PyTorch Hub.
\item \textbf{BarlowTwins} (non-contrastive SSL): Redundancy-reduction via cross-correlation matrix alignment \citep{Zbontar2021BarlowTwins}; official weights.
\end{itemize}

All backbones output 2048-dimensional features via global average pooling (\texttt{fc=Identity}). Backbones are frozen for all probe experiments.

\subsection{Spurious/Task Probe Accuracy Ratio}

For each frozen backbone, we train separate linear probes (logistic regression) to predict the \textit{spurious attribute} and the \textit{task label} from the 2048-dimensional features:

\begin{equation}
\text{ratio} = \frac{\text{spurious\_probe\_acc}}{\text{task\_probe\_acc}}
\end{equation}

A ratio above 1.0 indicates stronger spurious relative to task encoding. This scale-free metric decouples paradigm effects from absolute difficulty differences across datasets.

\subsection{Group-Balanced Probe Protocol}

Probes are trained on a \textit{group-balanced} subset of the validation split (equal examples per task label $\times$ spurious attribute group), not the training set. This avoids inflating spurious probe accuracy due to the 95\% Waterbirds training correlation. Evaluation uses the full balanced test split (Waterbirds: 5,794 images; CelebA: 720 images). Five random seeds per probe.

\subsection{Statistical Analysis}

Pairwise comparisons use two-sample $t$-tests across 5 seeds, Bonferroni-corrected for 6 pairs. Gate criterion: $p_\text{bonf} < 0.05$ and mean difference $\geq 0.02$. Cohen's $d$ uses pooled standard deviation (equal-$n$ two-sample formula: $d = (\mu_1 - \mu_2)/s_\text{pooled}$, $s_\text{pooled} = \sqrt{(s_1^2 + s_2^2)/2}$) to measure effect size. One-way ANOVA tests joint paradigm significance.

\subsection{Background-Replacement Augmentation (Mechanism Test)}

SimCLR trained from scratch on Waterbirds under two conditions: \textbf{SimCLR-Original} (standard augmentation) and \textbf{SimCLR-NoBackground} (adds random background replacement using Places365 images via CUB-200-2011 segmentation masks). Matched hyperparameters: SGD lr=0.03, batch=256, 50 epochs, NT-Xent temperature=0.5. Mechanism activation verified by pixel difference between views.

\subsection{Datasets}

\textbf{Waterbirds} \citep{Sagawa2020GroupDRO}: 4,795 train / 5,794 balanced test. Task = bird species; spurious = background type (95\% train correlation).

\textbf{CelebA} \citep{Liu2015CelebA}: 162,770 train / 720 balanced test. Task = Blond\_Hair; spurious = Male.
